\section{Canonical Definitions for This Draft}
\label{app:definitions}
The entries below consolidate the working definitions used in the body. They specify a model vocabulary, not a claim that each term already has an experimentally established physical referent.

\begin{description}[style=nextline,leftmargin=1.5em,font=\normalfont\bfseries]
\item[Quantum Relativity]
A theory of relational complexity governed by conservation of information. QR studies registered relational objects, their native complexity, and admissible information-preserving transductions. Quantum and geometric descriptions are investigated as compatible reductions of this common relational process.
\item[Thalean]
Named for Thales of Miletus and the search for explanation through natural principle rather than caprice. Water as an originating principle is historical background, not QR's literal substrate; see the naming note in the Introduction.
\item[Registered relational object]
A relational structure with a specified registration, retained distinctions, context, and equivalence of presentations. It is the typed object to which a complexity readout is applied.
\item[Relational complexity]
The nonnegative functional $\mathcal C(R)$ measuring how much relational organization is present in the declared registered object. It respects the chosen presentation equivalence. A general native formula is a construction target, not supplied by raw vertex count, path length, or terminology alone.
\item[Scalar]
A scalar readout of relational complexity, optionally encoded as $S(R)=G^{\mathcal C(R)}$ for a fixed dimensionless $G>1$. It is distinct from signed orientation, amplitude, and mechanical action.
\item[Complexity increment]
$\Delta\mathcal C_i=\mathcal C(R_{i+1})-\mathcal C(R_i)$ on one consistent accounting domain. It may be positive, zero, or negative; an event need not increase it.
\item[Scalar relay]
Reconstruction of scalar complexity from an admitted relational update and sufficient available context. Incremental relay obeys $S_{i+1}=S_iG^{\Delta\mathcal C_i}$ and needs a known starting value or equivalent reference. A scalar need not be sent as an independent payload, but its informational prerequisites cannot be omitted.
\item[Thalean conservation of information]
The complete-output recovery requirement for admitted transduction. In the finite classical description it is injectivity; in the stated coherent representation it is isometry. It does not identify complexity with entropy or assert an additive sum of independent scalar and phase information.
\item[Thalion]
The structured relational transducer carried by $G_{60}\cong\mathrm{AT4val}[60,6]$, together with only the quotient, deck, transport, closure, history, and readout structures explicitly supplied or derived in the construction under discussion.
The Tolkien resonance ``steadfast'' is a secondary mnemonic for informational fidelity, not the source of this mathematical definition.
\item[Thalean transducer]
A Thalion considered through the admissible transformations it performs and the distinctions it propagates, retains, compares, or exposes under projection. Apparent erasure in a projection must not be confused with erasure from a complete admissible event.
\item[Relation]
A locally available relational possibility, represented where appropriate by adjacency or a typed interaction interface. Availability is not itself a realized transmission.
\item[Continuation]
An admissible possible next relation from a specified current configuration and registration context.
\item[Informative action]
An admitted realized continuation that creates or changes a registered relational distinction. This primitive term is not identified with a dimensionful mechanical action functional.
\item[Admissible transduction]
A local transformation satisfying the declared relational constraints and preserving the information required for reconstruction from its complete output. In the classical finite formulation it is injective on its admitted complete-input domain.
\item[Propagated component]
The output eligible to acquire an adjacent field address under the admitted handoff rule. It need not alone encode the complete input.
\item[Trace]
A retained relational consequence of an informative action.
\item[Receipt]
A retained complementary distinction required, together with the propagated component and relevant memory, to account for a completed admissible handoff. A receipt can be trivial when the propagated component is already sufficient.
\item[Relational Spin Residue]
A name reserved for a specifically established spin-related central or projective residue. It is narrower than the generic receipt and does not confer a physical spin interpretation on arbitrary finite holonomy.
\item[QR event]
In the propagation sector, one completed, receipted, admissible relational handoff: local transduction, an adjacent address acquisition, and retained complementary accounting. Internal registration actions are not automatically field displacements.
\item[Event primacy]
The methodological principle that becoming is modeled through completed relational occurrences and their composition, with states representing conditions before, after, or across those occurrences.
\item[Becoming]
The ordered composition of admitted realized relational events through which relational organization is transformed with complete informational accounting preserved. Complexity readouts describe aspects of this process; they do not define the occurrence order.
\item[Registered history]
The compositional record of realized events, including retained distinctions. It is richer than a list of visible endpoint configurations.

\item[Mode]
The $S_3$ orientation parity of the registered event cycle. On the declared free closure-frame orbit with reference $f_*$, $\mu(g\cdot f_*)=\sgn(g)$. The generator $\Cevt$ has order three and the register-reversing reflection $\Mevt$ has order two, with $\Mevt\Cevt\Mevt=\Cevt^{-1}$. The two orientation classes are intrinsic; their positive and negative names depend on the reference.
\item[Registered event phase]
The $C_3$ position $p$ in a registered modal-phase frame $\Cevt^p\Mevt^m\cdot f_*$. Phase and Mode are semidirectly related by the $S_3$ action, not commuting independent dynamics.
\item[Event-presentation invariant]
The declared event type or relational structure preserved across its admitted presentation orbit. This does not identify distinct completed occurrences or replace the propagation-event definition.
\item[Collective eigenmode]
An eigenvector or invariant subspace of a specified global update operator, distinct from registered-event Mode.
\item[Delta]
$\delta_e$ denotes a registered distinction of an event; $\Delta[\Gamma]$ denotes a specified retained composition of such distinctions. Neither notation is automatically a scalar duration or a universal additive quantity.
\item[Complete operational state]
The current configuration together with retained distinctions relevant to future admitted response. This operational notion is distinct from a complete microscopic information-preserving description.
\item[Predictive history class]
An equivalence class of histories having the same admissible future language and the same registered responses to each future test. Equality only on a possibly empty intersection of future languages is insufficient.
\item[Memory]
Retained distinction with future discriminating power in the specified apparatus or transducer.
\item[Closure]
An admitted operation bringing relational histories into a common comparison context. It does not necessarily erase their difference or restore their complete initial state.
\item[Closure residue]
A retained distinction that survives an admitted comparison or closure.
\item[Closure feedback]
A proved dependence of a later admissible continuation on a previously retained closure residue. The existence of a residue alone does not establish feedback.
\item[Thalean time]
The primitive additive oriented winding of the unique $G_{30}$-visible registered-history cycle $r_1$ in the certified Project41 cycle register. Its immediate domain is the admitted closed registered histories.
\item[Thalean tick]
One positively oriented primitive traversal of $r_1$, normalized by $\Delta\tth=+1$. It is not generally one primitive event step or one physical second.
\item[Event depth]
Irreducible primitive composition depth in a specified admitted grammar. It is nonnegative and must not be reduced by renaming a multi-step process as one macro-event.
\item[Event-preserving history]
An admitted composition that preserves the declared relational structure being propagated through its transductions and handoffs, including the complementary accounting required for reconstruction.
\item[Thalean causal datum]
$\cth=\sup_\gamma d_{\mathrm{field}}(s(\gamma),t(\gamma))/|\gamma|_{\mathrm{event}}=1$ in the native local normalization, over positive-depth admitted event-preserving histories with a saturating handoff available.
\item[Transport cocycle]
A displacement assignment compatible with history composition. It is additive in a common abelian register, or obeys a specified twisted composition law when frames change.
\item[Measure]
A weighting of alternatives supplied by a declared relational and instrument construction.
\item[Registered measure]
The realization of that weighting through an admitted output transducer, preserving the intended observable rather than silently replacing its measure.
\item[Frequency realization]
A registered sequence whose conditional outcome counts realize or converge to the specified measure.
\item[Process provenance]
The admitted native law generating the chronological sequence, distinct from a proof that some sequence can realize the measure.
\item[Registration nonamplification]
A common downstream registration cannot distinguish identical upstream outputs without an additional input, changed instrument, or other new resource.
\item[Local orientation cell]
A four-state free $V_4$ fiber with three nonidentity pairing operations. These pairings are the perfect matchings of an abstract $K_4$, not a claim that its six edges are carrier adjacencies.
\item[Carrier triality]
The abstract $S_3$ action permuting the three nonidentity elements of $V_4$. Its extension to a particular carrier action or execution grammar requires a separate check.
\item[Sheet]
A declared relational, spectral, quotient, or cover presentation. The word does not imply a physical surface or spatial dimension.
\item[Continuation multiplicity]
The number, weight, or dimension of admitted futures over a stated event depth and under a stated equivalence relation.
\item[Continuation entropy]
A specified measure of growth or diversity of admitted futures. A constant covering multiplicity alone does not change an asymptotic exponential growth rate.
\item[Frequency]
Recurrence or eigenphase per a declared evolution parameter. A physical angular frequency additionally requires calibration of that parameter.
\item[Temperature]
An ensemble parameter associated with an appropriate thermal occupation description, after the energy and statistical assumptions have been supplied. Generic eigenmode occupation is not necessarily thermal.
\item[Balanced release]
Efficient local conversion of internal excitation into admitted outgoing eigenmodes with complete information accounting retained. It is not defined as maximal excitation or equality of transmitted and retained scalar quantities.
\end{description}

\subsection*{Additional Draft 3 terminology}
\paragraph{Common cover.} $\mathcal K_{120}=\mathrm{C4}[120,38]$, with specified projections to the three census children. Its extra coordinate over $G_{60}$ is canonical-double sheet parity.
\paragraph{Shared family intermediate.} $Q_{30}=L(\mathrm{CDC}(P))$, distinct from the inherited $G_{30}=L(D)$.
\paragraph{Relational kernel.} The eight-relation coherent configuration $\mathscr K_{\rm rel}$ on $X_{12}=C_3\times V_4$ defined in Section~\ref{sec:relkernel}. It is neither the common-cover graph nor the history conference operator $K_{\rm hist}$.
\paragraph{Informative coupling.} Dependence of admitted continuation or registered response on a retained distinction, with the comparison context and other retained variables fixed.
\paragraph{Information current.} Transport of an operationally distinguishable payload through admitted field handoffs and a specified downstream decoder.
\paragraph{Immediate alternatives.} Distinct representatives of a declared registration fiber. A fixed cover gives no additional per-step branching once the initial lift is chosen; amplitudes and a quantum instrument are additional structures.
\paragraph{Frame and history actions.} Frame isotropy fixes a selected graph position or presentation. A history action is defined on a succession or representation domain. Equal abstract group names do not identify those actions.
